% !Mode:: "TeX:UTF-8"
\chapter{相关技术与理论基础}

\section{深度学习框架与模糊测试技术}

\subsection{深度学习框架的体系结构}

当前主流DL框架普遍采用分层架构设计\cite{ref_pytorch_arch}。
以PyTorch为例，其体系结构自顶向下可划分为三个层次。
最上层是面向用户的Python前端接口，提供张量操作、自动微分以及神经网络模块的高级抽象。
中间层为计算图的中间表示（Intermediate Representation），
负责将用户定义的运算逻辑转化为有向无环图（Directed Acyclic Graph, DAG）结构，
其中节点对应算子调用，有向边描述数据的流向依赖关系。
最底层是C++实现的ATen张量计算库与CUDA核函数调度引擎，
直接操控GPU显存分配、并行线程调度与底层数值运算。

这种三层结构的设计使得框架接口简洁、执行效率高。
但同时也意味着，当底层ATen库中某个算子存在边界条件处理缺陷时，
错误会被上层抽象层层封装，
上层用户极难通过Python级别的堆栈信息追溯至真实的缺陷触发位置。
本课题所构建的故障溯源引擎正是针对这一跨层封装带来的定位困难而设计。

\subsection{模糊测试技术原理与分类}

Fuzzing的核心思想是通过自动化手段持续生成大量测试输入，
将其送入待测程序执行，并监控运行时的异常信号\cite{ref_fuzzing_art}。
依据输入的构造策略，Fuzzing可分为两大类别：
基于变异（Mutation-based）的方法从已有的种子输入出发，
通过随机翻转、插入或删除字节来派生新的测试用例\cite{ref_afl}；
基于生成（Generation-based）的方法则依据输入格式的语法规范或语义模型
从零构建符合特定约束的测试输入\cite{ref_grammar_fuzzing}。

面向DL框架的Fuzzing面临两项特有挑战。
其一，输入空间维度远高于传统软件：
测试用例不再是单一的二进制文件，而是包含多维张量属性、算子类型枚举及其拓扑组合的复合结构。
其二，"崩溃"的定义更为复杂：
除段错误与未捕获异常外，数值精度的静默偏离同样构成框架缺陷，
但此类异常不会触发操作系统信号，需要额外的差分预言机制进行检测\cite{ref_differential_testing}。
本课题的测试管理系统需要对上述异构崩溃产物统一归档并进行特征提取，
这对底层数据流水线的泛化能力提出了明确要求。

\section{文本语义表征与向量检索技术}

\subsection{基于孪生网络的预训练语义编码模型}

传统的词袋模型或TF-IDF方法将文本映射为稀疏高维向量，
无法捕捉词序与上下文语义信息\cite{ref_tfidf}。
Reimers等人提出的Sentence-BERT\cite{ref_sbert}
在预训练语言模型的基础上引入孪生网络（Siamese Network）结构，
通过对比学习目标使模型输出的句子级嵌入向量在语义空间中具备距离可度量性。
具体而言，两条输入文本分别经过共享参数的编码器后，
各自产出一个固定维度的稠密浮点向量，
余弦相似度即可直接反映两段文本的语义接近程度。

本课题选用all-MiniLM-L6-v2\cite{ref_minilm}作为具体的编码器实例。
该模型基于6层Transformer架构蒸馏而来，
输出维度为384维，在保持较高语义区分度的同时将推理延迟控制在毫秒级。
在本系统中，该模型被部署于数据摄入流水线的向量化环节，
负责将崩溃日志的报错堆栈文本编码为稠密语义向量，
为后续的相似缺陷检索与聚类分析提供数值基础。

\subsection{高维向量近似最近邻检索}

当向量集合的规模增长至百万级别时，
逐一计算查询向量与所有库中向量之间的距离将带来不可接受的延迟。
ANN检索通过牺牲少量精度来换取数量级的速度提升\cite{ref_ann_benchmark}。
Milvus\cite{ref_milvus_arch}是一款面向海量向量数据的开源检索引擎，
支持多种索引类型与距离度量方式。

本课题选用IVF\_FLAT（Inverted File Index with Flat quantization）作为主索引结构。
该索引首先通过聚类算法将向量空间划分为若干Voronoi单元，
查询时仅在距离查询点最近的若干单元内执行精确扫描，
从而将搜索范围从全量集合缩减至局部子集\cite{ref_ivf}。

向量间的距离度量采用余弦相似度。对于两个$d$维向量$\mathbf{a}$与$\mathbf{b}$，
其余弦相似度定义为：
\begin{equation}
\label{eq:cosine}
\text{sim}(\mathbf{a}, \mathbf{b}) = \frac{\sum_{i=1}^{d} a_i b_i}{\sqrt{\sum_{i=1}^{d} a_i^2} \cdot \sqrt{\sum_{i=1}^{d} b_i^2}}
\end{equation}
其中$a_i$与$b_i$分别为向量$\mathbf{a}$和$\mathbf{b}$的第$i$个分量。
该度量对向量的模长不敏感，仅关注方向一致性，
适用于本系统中长度不等的崩溃文本经编码后产生的语义向量之间的相似度计算。

\section{降维分析与不变量推断方法}

\subsection{主成分分析原理}

PCA是一种经典的线性降维方法\cite{ref_pca}，
其目标是将高维数据投射至低维子空间的同时保留最大方差信息。
给定$n$个$d$维样本构成的数据矩阵$\mathbf{X} \in \mathbb{R}^{n \times d}$（已中心化），
PCA首先计算样本协方差矩阵：
\begin{equation}
\label{eq:cov}
\mathbf{C} = \frac{1}{n-1}\mathbf{X}^{\top}\mathbf{X}
\end{equation}
其中$\mathbf{C} \in \mathbb{R}^{d \times d}$为对称半正定矩阵。
随后对$\mathbf{C}$执行特征值分解：
\begin{equation}
\label{eq:eigen}
\mathbf{C}\mathbf{v}_k = \lambda_k \mathbf{v}_k, \quad k = 1, 2, \ldots, d
\end{equation}
其中$\lambda_k$为第$k$个特征值，$\mathbf{v}_k$为对应的单位特征向量。
将特征值按降序排列后，取前$p$个特征向量构成投影矩阵$\mathbf{W} = [\mathbf{v}_1, \ldots, \mathbf{v}_p] \in \mathbb{R}^{d \times p}$。
原始数据通过$\mathbf{Y} = \mathbf{X}\mathbf{W}$即可降阶至$p$维表示。

本课题将384维的测试用例语义向量通过PCA降至2维平面坐标，
用于生成测试集多样性的空间分布散点图。
该降维结果为研究人员提供了对测试集探索广度的直观监测手段。

\subsection{Daikon动态不变量检测框架}

Daikon\cite{ref_daikon}是由Ernst等人提出的动态不变量检测工具。
其核心思想是：从程序多次执行的运行时轨迹中，
自动推断可能成立的程序不变量（Program Invariant）。
具体而言，Daikon在被观测程序的关键执行点采集变量的实际取值，
然后将这些取值序列与预定义的不变量模板库（如$x > 0$、$x = a \cdot y + b$、$x \in \{v_1, v_2, \ldots\}$等）进行匹配。
若某一模板在全部已观测轨迹上均成立，则将其报告为一条候选不变量。

本课题借鉴了Daikon的核心推断范式，但应用场景与经典用法存在本质差异。
经典Daikon面向的是单一程序的正常执行轨迹，
而本系统面对的是DL框架Fuzzing产物中"成功执行"与"崩溃触发"两组对照样本的参数空间。
系统从中推断出区分两组样本的数学边界条件，
将其作为算子输入参数的物理约束红线。
这一改造的具体算法设计将在第四章详细阐述。

\section{检索增强生成技术}

RAG架构\cite{ref_rag_survey}由检索器（Retriever）与生成器（Generator）两个核心组件构成。
当用户发起查询时，检索器首先从外部知识库中召回与查询语义相关的文档片段，
随后将这些片段作为上下文拼接至LLM的输入提示词中，
引导生成器在受约束的知识范围内产出回答。

这一架构相较于直接调用LLM具有两方面优势。
其一，通过外部知识注入减轻模型的参数记忆负担，降低Hallucination的发生概率。
其二，知识库可动态更新而无需重新训练模型权重，
使系统能够适应持续增长的领域数据\cite{ref_rag_knowledge}。

本课题在两个场景中集成RAG机制：
一是故障辅助诊断——以当前崩溃用例的语义向量召回历史相似案例，
为LLM提供具体的项目环境上下文；
二是算子缺口补盲——将识别出的未覆盖算子组合信息注入提示词，
驱动LLM生成具有针对性的测试脚本。
RAG在本系统中的具体集成方式将在第四章展开。

\section{代码覆盖率采集技术}

代码覆盖率是衡量测试完备程度的基础指标\cite{ref_coverage_survey}。
GNU工具链中的Gcov\cite{ref_gcov}是面向C/C++程序的覆盖率采集工具。
其工作原理分为三个阶段：
在编译阶段，通过向GCC传递\texttt{--coverage}选项，
编译器在源代码的每个基本块入口处插桩（Instrumentation）计数器指令，
同时生成记录源码结构信息的\texttt{.gcno}文件；
在执行阶段，程序运行时各计数器累加实际命中次数，
进程正常退出后将计数结果落盘为\texttt{.gcda}文件；
在分析阶段，Gcov读取\texttt{.gcno}与\texttt{.gcda}文件，
计算并输出每个源文件的行级与分支级覆盖率。

LCOV\cite{ref_lcov}是对Gcov输出的上层封装工具。
它支持对多次运行产生的覆盖率数据执行合并（Merge）与差集运算，
并可生成HTML格式的可视化覆盖率报告。
在本课题中，系统需要在插桩编译的PyTorch底层C++引擎上
逐条执行测试用例并采集对应的覆盖率增量数据。
LCOV的合并功能使系统能够追踪每条用例对全局覆盖率的边际贡献，
这一数据是计算CEI的必要输入。

\section{本章小结}

本章介绍了本课题涉及的核心技术与理论基础。
首先阐述了DL框架的分层体系结构与Fuzzing技术的基本原理，明确了待测对象与测试方法的技术特征。
随后介绍了Sentence-BERT语义编码模型与Milvus向量检索引擎的工作机制，为系统的语义检索能力提供了技术支撑。
在此基础上，说明了PCA降维方法与Daikon不变量推断的核心思想，为后续的多样性监测与边界定界奠定了理论依据。
最后概述了RAG架构与Gcov/LCOV覆盖率采集工具链，分别对应系统的智能诊断模块与效能评估模块的底层技术需求。
上述技术将在第四章中结合系统的具体算法设计进行深入阐释。
